\section{Open problems}\label{sec:open}

The arbitrary-field L\"uck approximation conjecture is refuted by Theorem~A, which is formal except for Theorem~H. The
following questions remain.

\begin{enumerate}[label=\arabic*.,leftmargin=2em]
\item \emph{A formal or elementary proof of Theorem~H.} Formalize the geometry of \cref{sec:hypothesis}, or find another
proof of (KN), for this or any square matrix of full inner rank over a free group in positive characteristic. For the
matrix~$A$ such a proof cannot apply to all covers of the base graph, and it must fail in characteristic~$0$
(\cref{rem:congruence}); heuristically, it must use the arithmetic of congruence covers.

\item \emph{Characteristic~$2$ and other primes.} Does part~(2) fail in every positive characteristic? The case $p=2$
concerns \emph{Bounded Nesting and Depth Reduction for Matrix Transposition on Tapes}, by the same authors. That paper asks
whether binary codes lifted along normal towers of covers with trivial intersection can have positive asymptotic rate; it
does not cover $2$-adic analytic towers, for example, and notes that a positive limit there would contradict the
characteristic-$2$ case of the conjecture~\cite[Section~8.1]{BoundedNesting2026}. Theorem~A shows that the conjecture
fails in characteristic~$3$, so the general conjecture no longer supports its characteristic-$2$ case. For $2$-adic
analytic towers, whose completions are virtually pro-$2$, Jaikin-Zapirain's refined conjecture~\cite{JaikinBlog2024}
still predicts the universal limit, and Theorem~A does not touch it. Goren--Oort configurations at other inert primes, where the supersingular locus is again a
union of projective lines indexed by a Bruhat--Tits tree, are natural candidates in other characteristics.

\item \emph{Part~(1) along the congruence chain.} Does $\rk_{F/N_k}(A)$ converge? If
$\dim_{\F_9}T_{\mathrm{tw}}(X_k)=V_k+c_k-2$ for every~$k$, which holds when $h^1(\OO_{S_k})=0$ in characteristic~$3$
(\cref{sec:RR}), then the normalized kernel over $\F_3$ is $22+2(c_k-2)/11^{3k}$, which tends to $22$ because the number
of cusps $c_k$ grows only like $11^k$; the ranks would then converge to~$88$. The formula holds at $k=0$, where the kernel is certified
(\cref{rem:base}), and, by computation, at $k=1$ (\cref{rem:numerics}).

\item \emph{The refined conjecture.} Jaikin-Zapirain's refined conjecture for chains whose completion is virtually
pro-$p$~\cite{JaikinBlog2024} is untouched: the bad chain here is pro-$11$ and the matrix lives in characteristic~$3$.

\item \emph{Other free groups.} Theorem~A concerns the free group of rank~$89$. Does part~(2) fail in characteristic~$3$
for every non-abelian free group of finite rank, for instance for the free group of rank~$2$, which contains the free
group of rank~$89$ as a normal subgroup of index~$88$?
\end{enumerate}
